%! Logarithmic Function Manipulation — Unit 5, Lesson 5.4 — Lesson Plan
\documentclass[10pt]{article}
\usepackage{precalculus-article}
\usepackage{precalculus-boxes}
\usepackage{graphicx}
\graphicspath{{images/}}
\newcommand{\TallMath}[1]{$\displaystyle #1\rule[-1.4em]{0pt}{3.2em}$}
\newcommand{\CourseName}{Precalculus}
\newcommand{\MeetingLength}{60 minutes}
\newcommand{\UnitNumberName}{Unit 5: Logarithmic Functions \quad}
\newcommand{\LessonNumberName}{Lesson 5.4: Logarithmic Function Manipulation}

\begin{document}
\begin{center}
    {\Large\bfseries \CourseName \\
    {\small\bfseries \UnitNumberName \LessonNumberName}}
\end{center}

\begin{tcolorbox}[colback=sky, colframe=navy, boxrule=0.9pt, arc=2mm,
    left=2mm, right=2mm, top=1.5mm, bottom=1.5mm]
    \textbf{Primary Objective:} Students rewrite logarithmic expressions using the product,
    power, and change-of-base properties and connect these algebraic manipulations to graphical
    transformations (vertical translations and dilations). Students also introduce the natural
    logarithm $\ln(x) = \log_e(x)$.
    \hfill\textit{\small Skills: AP 1.B, AP 3.A \quad LO 2.12.A}
\end{tcolorbox}

\renewcommand{\arraystretch}{1.6}
\begin{skillbox}[Priority Ideas \& Skills]{goldbox}
\begin{tabularx}{\linewidth}{@{}X X@{}}
\textbf{AP Skills} &
\textbf{Key Understandings} \\
\midrule
\textbf{AP 1.B} --- Express equivalent forms analytically. &
\begin{itemize}[leftmargin=1.2em,itemsep=2pt,topsep=0pt]
  \item Product property: $\log_b(xy) = \log_b x + \log_b y$; graphically, every
        horizontal dilation $\log_b(kx)$ is a vertical translation $\log_b k + \log_b x$.
        (EK 2.12.A.1)
  \item Power property: $\log_b(x^n) = n\log_b x$; graphically, $\log_b(x^k)$ is a
        vertical dilation of $\log_b x$. (EK 2.12.A.2)
\end{itemize} \\
\textbf{AP 3.A} --- Describe characteristics of a function. &
\begin{itemize}[leftmargin=1.2em,itemsep=2pt,topsep=0pt]
  \item Change of base: $\log_b x = \tfrac{\log_a x}{\log_a b}$, so all logarithmic
        functions are vertical dilations of each other. (EK 2.12.A.3)
  \item $\ln(x) = \log_e(x)$; the natural logarithm with base $e \approx 2.718$.
        (EK 2.12.A.4)
\end{itemize} \\
\end{tabularx}
\end{skillbox}

\begin{skillbox}[{Vocabulary, Concepts, \& Theorems}]{greenbox}
\renewcommand{\arraystretch}{1.4}
\begin{tabularx}{\linewidth}{@{} >{\leavevmode\bfseries\color{navy}}p{0.30\linewidth} X @{}}
Product property for logarithms & $\log_b(xy) = \log_b x + \log_b y$; a horizontal dilation
                                  is equivalent to a vertical translation of the graph. \\
Power property for logarithms   & $\log_b(x^n) = n\log_b x$; raising the input to a power
                                  is equivalent to a vertical dilation of the graph. \\
Change of base property         & $\log_b x = \dfrac{\log_a x}{\log_a b}$ for any valid base
                                  $a$; allows evaluation on any calculator. \\
Natural logarithm ($\ln$)       & $\ln(x) = \log_e(x)$, the logarithm with base
                                  $e \approx 2.718$. \\
Vertical translation            & A shift of the graph up or down; produced by adding/subtracting
                                  a constant outside a function. \\
Vertical dilation               & A stretch or compression of the graph toward/away from the
                                  $x$-axis; produced by multiplying the output by a constant. \\
\end{tabularx}
\end{skillbox}

\begin{skillbox}[Activate Prior Knowledge \& Spiral Review]{sky}
    Spiral review (warm-up) covers: evaluating basic logarithm expressions ($\log_2 32$,
    $\log_3 81$, $\log_{10} 1000$); identifying the additive pattern in $f(x) = \log_5 x$
    as the input multiplies by 5; and making an analogy between the exponent rule
    $2^3 \cdot 2^4 = 2^7$ and a predicted sum rule for logarithms. These directly activate
    EK 2.12.A.1 and prepare students for the product property.
\end{skillbox}

\begin{skillbox}[Hook]{sky}
    Write on the board: \textbf{We know $\log_2 8 = 3$ and $\log_2 4 = 2$.}

    Ask: ``What do you predict $\log_2(8 \cdot 4) = \log_2(32)$ equals? Can you write it as
    a sum of the two values above?'' Give students 60 seconds to discuss with a partner.
    Most will predict $3 + 2 = 5$. Verify: $\log_2(32) = 5$ because $2^5 = 32$. Then ask:
    \textbf{``Is there a general rule? When can we break a log of a product into a sum?''}
    This motivates the product property.
\end{skillbox}

\begin{skillbox}[Lesson]{sky}
\begin{multicols}{2}
\textbf{Part 1 (12 min): Product Property (EK 2.12.A.1).}

State: $\log_b(xy) = \log_b x + \log_b y$. Verify with the hook example. Then reveal the
graphical connection: $f(x) = \log_b(kx) = \log_b k + \log_b x$, so a horizontal dilation
by factor $k$ is the same graph as a vertical translation up by $\log_b k$. Have students
complete the product-property fill-in in the notes.

\columnbreak

\textbf{Part 2 (10 min): Power Property (EK 2.12.A.2).}

State: $\log_b(x^n) = n\log_b x$. Show that $f(x) = \log_b(x^k) = k\log_b x$, i.e., a
vertical dilation by $k$. Work through $\log_3(9^2)$ together. Students complete the
power-property fill-in.

\textbf{Part 3 (8 min): Change of Base (EK 2.12.A.3).}

State: $\log_b x = \tfrac{\log x}{\log b}$ (common log form). Implication: every
$\log_b$ is a constant vertical dilation of $\log_{10}$. Evaluate $\log_5 25$ on a
calculator. Then introduce $\ln(x) = \log_e(x)$; evaluate $\ln(e)$ and $\ln(e^3)$
(EK 2.12.A.4).
\end{multicols}
\end{skillbox}

\begin{skillbox}[Active Monitoring]{redbox}
\begin{itemize}[leftmargin=1.2em,itemsep=2pt,topsep=0pt]
  \item After Part 1: cold-call ``What does the product property tell us about the graph of
        $y = \log_2(8x)$ compared to $y = \log_2 x$?'' Expect: vertical translation up by 3.
  \item Watch for students who write $\log_b(x + y) = \log_b x + \log_b y$ (addition, not
        multiplication inside) --- address this confounding error immediately.
  \item After Part 2: ask ``Is $\log_3(x^4) = [\log_3 x]^4$?'' Expect: No;
        $\log_3(x^4) = 4\log_3 x$. Students confusing the exponent rule with power rule are
        common.
  \item After Part 3: confirm students can evaluate $\log_5 17$ on their calculator using
        $\log(17)/\log(5)$ and that they understand this is exact (not approximate).
  \item Check that students correctly identify $\ln(e) = 1$ without a calculator.
\end{itemize}
\end{skillbox}

\begin{skillbox}[Group Work \& Differentiation]{redbox}
\begin{itemize}[leftmargin=1.2em,itemsep=2pt,topsep=0pt]
  \item \textbf{Tier R (Remediate):} Apply product and power properties to straightforward
        numerical expressions: $\log_2(4 \cdot 8)$, $\log_3(9 \cdot 27)$, $\log_5(25 \cdot 125)$;
        then $\log_2(4^2)$, $\log_3(27^3)$.
  \item \textbf{Tier A (Apply):} Rewrite $\log_2(32x)$ as $a + \log_2 x$; rewrite $\log_3(x^4)$
        as a vertical dilation; evaluate $\log_5(17)$ using change of base; show
        $\log_2(x) = \ln(x)/\ln(2)$.
  \item \textbf{Tier E (Extension):} Prove $\log_3(9x) = \log_3 x + 2$ and identify the
        transformation; prove all $\log_b$ are vertical dilations of $\log_{10}$; simplify
        $\log_2(8x^3)$ using both properties with full justification.
\end{itemize}
\end{skillbox}

\begin{skillbox}[Individual Work \& Assessment]{redbox}
Exit ticket (2 questions, 1 page):
\begin{enumerate}[leftmargin=1.5em,itemsep=2pt,topsep=2pt]
  \item Rewrite $\log_4(16x^2)$ in the form $a + b\log_4 x$; identify $a$ and $b$.
        (Tests product + power properties together, EK 2.12.A.1--A.2.)
  \item Use change of base to express $\log_7(50)$ as a ratio of common logarithms.
        (Tests EK 2.12.A.3.)
\end{enumerate}

Sort by result: students missing \#1 need product/power practice; students missing \#2
need change-of-base practice and/or calculator fluency.
\end{skillbox}

\begin{skillbox}[Reinforcement \& Extension]{goldbox}
\textbf{Homework (6 problems):} product property numerical \& algebraic; power property
numerical \& algebraic; graphical connection for $f(x) = \log_2(4x)$; change-of-base
explanation and dilation reasoning; evaluating natural logarithms ($\ln(e^5)$, $\ln(1)$,
$\ln(e^{-2})$); AP-style MC on equivalent forms.

\textbf{Extension:} Prove the quotient rule $\log_b(x/y) = \log_b x - \log_b y$ using the
product property and the property $b^{-n} = 1/b^n$.

\textbf{Preview of Lesson 2.13:} Solving exponential and logarithmic equations --- using
properties developed today to isolate variables, including equations with $\ln$.
\end{skillbox}

\end{document}
